\section{Introduction}\label{dag:sec:intro}

The two preceding chapters, for all their differences, share an
assumption so ingrained it is easy to miss: the unit of optimization
is one SQL statement. Search-based rewriting perfected it for queries
within its reach, and LLM-based rewriting extended it to queries
beyond; neither ever looks past the semicolon. Modern analytics
increasingly does. SQL today is less a language for asking individual
questions than an implementation language for recurring data products,
governed dashboards, AI-ready feature tables, self-service BI
datasets, and organization-wide reporting layers, and much of it is
generated rather than written from scratch: recent studies report that
70\% of analytics professionals use AI for code
development~\cite{dbt2025state}, and that BI-generated SQL contains
extremely complex scalar expressions and relational operator
trees~\cite{getreal2018}. The result is not just harder queries but a
different shape of workload, and it is the shape, not the queries,
that this chapter optimizes.

\ph{The Rise of SQL Pipelines}
That shape is the SQL pipeline. Data teams increasingly encode
business logic as graphs of transformations that feed one another on a
schedule, driven by logic that must be computed repeatedly and reused
consistently across dashboards, feature tables, compliance reports,
and downstream applications. One of the most widely adopted tools for
these workflows is \emph{dbt} (the data build
tool)~\cite{dbtWhatExactly}, reported to be used by more than 50,000
data teams weekly~\cite{dbtComplexity}; it lets teams write
transformations as version-controlled SQL \code{SELECT} statements
(called ``models'' in dbt, a term we use interchangeably in this
chapter), declare the dependencies among them, and compile the
resulting graph into SQL jobs executed by a data warehouse.
\Cref{dag:fig:intro-toy} gives a toy example: a raw orders table feeds
a cleaning transformation whose output feeds both a revenue table and
a daily active-customer table; the SQL text explains each local
transformation, while the edges explain how the transformations
compose. The pattern is much broader than dbt:
SQLMesh~\cite{sqlmeshOverview}, Dagster asset
graphs~\cite{dagsterAssets}, Airflow directed acyclic
graphs~\cite{airflowDags}, Prefect flows~\cite{prefectTasks}, Argo
Workflows~\cite{argoDag}, Databricks Workflows~\cite{databricksJobs},
and medallion lakehouse pipelines~\cite{databricksMedallion} all
represent recurring transformations through explicit dependencies. We
target the SQL subset of this broader pattern, which we call a
\emph{SQL pipeline DAG}: a recurring directed acyclic graph whose
nodes are SQL transformations or SQL-produced tables and views, and
whose edges record that one node reads another node's output.

\begin{figure}[t]
\centering
\begin{tikzpicture}[
  font=\scriptsize,
  node distance=5mm and 4mm,
  box/.style={draw=black!70, text=black, rounded corners=2pt, align=left,
              inner sep=2.2pt, minimum width=22mm},
  edge/.style={-{Stealth[length=1.5mm]}, black!70, thin}
]
\node[box, minimum width=17mm] (orders) {\textbf{orders}\\raw table};
\node[box, minimum width=24mm, right=of orders] (clean) {\textbf{clean\_orders}\\\code{SELECT ...}\\\code{FROM orders}\\\code{WHERE status='paid'}};
\node[box, minimum width=20mm, above right=of clean] (rev) {\textbf{daily\_revenue}\\\code{GROUP BY day}};
\node[box, minimum width=20mm, below right=of clean] (cust) {\textbf{active\_customers}\\\code{SELECT DISTINCT}\\\code{customer\_id}};
\draw[edge] (orders) -- (clean);
\draw[edge] (clean) -- (rev);
\draw[edge] (clean) -- (cust);
\end{tikzpicture}
% \inv
\caption{A toy SQL pipeline DAG. Nodes are SQL transformations or SQL-produced tables (called ``models'' in dbt); edges show which transformation consumes another transformation's output.}
\label{dag:fig:intro-toy}
\inv
\end{figure}

\ph{The Pipeline Optimization Gap}
Pipelines of hundreds or thousands of interdependent transformations
sit in a gap no existing technique covers. Single-statement rewriting,
whether rule-based~\cite{wetune}, synthesis-based
(\cref{slab:chap})~\cite{dong2023slabcity}, or LLM-based
(\cref{gen:chap})~\cite{genrewrite,sun2024r}, cannot see opportunities
whose correctness or benefit depends on the surrounding pipeline: work
no final output consumes, equivalent logic duplicated in distant
transformations, refresh schedules mismatched to how inputs change.
Inlining every referenced transformation into each downstream query,
to give a single-query rewriter global view, produces SQL nested tens
or hundreds of layers deep, well beyond what today's optimizers and
rewriters can even tackle. Classical multi-query
optimization~\cite{sellis1988multiple,roy2000efficient} considers
multiple queries together, but its goal is to share common scans,
joins, or intermediate results across a submitted batch; it does not
rewrite a pipeline by adding, removing, splitting, or merging
transformations based on downstream demand and refresh behavior.
Materialized-view
selection~\cite{harinarayan1996implementing,agrawal2000automated}
precomputes intermediate results that have already been identified,
but does not by itself determine that a pipeline should drop unused
work, factor repeated logic written in different forms, move expensive
operations after data reduction, or refresh different regions at
different rates. And workflow
schedulers~\cite{airflowDags,dagsterAssets,prefectTasks,argoDag} know
the dependencies but optimize execution order, failure recovery,
retries, and resource allocation, never the SQL semantics of the
transformations themselves. The pipeline as a semantic object has no
optimizer.

\ph{Dependencies as Signals}
\label{dag:sec:dependencies-as-signals}
This chapter's starting observation, introduced in Chapter~1, is that
explicit dependencies are not merely execution constraints; they are
signals that can be leveraged for optimization. A SQL pipeline DAG
exposes six kinds of information that are hidden, or much harder to
recover, when queries are optimized independently: producer-consumer
edges reveal how one transformation's output is used by later
transformations (\textbf{edge semantics}); similar business logic can
appear in distant parts of the graph even when the SQL text differs
(\textbf{cross-graph redundancy}); downstream dependencies reveal
which columns, joins, filters, and intermediate results can affect
final outputs (\textbf{downstream demand}); the graph reveals where
expensive operations occur relative to data-reducing filters,
projections, joins, and aggregations (\textbf{operator placement});
unlike single-query rewriting, which is constrained to equivalent
rewrites that are cheaper in isolation, pipeline rewriting can
consider alternatives that produce additional reusable outputs and are
locally more expensive, as long as those outputs are materialized once
and reused by enough downstream consumers
(\textbf{rewrite-persistence coupling}); and schedules and source
update rates reveal how often inputs change and outputs are consumed
(\textbf{refresh patterns}).

\ph{Our Approach}
\dagsmith, the system this chapter presents, is to the best of our
knowledge the first holistic, dependency-aware rewriting system for
SQL pipeline DAGs. It converts the six signals into six classes of
optimization: \textbf{(1) dependency-edge simplification:} remove or
simplify dependency edges whose purpose can be expressed more
directly, while preserving final outputs; \textbf{(2) non-local
semantic reuse:} factor repeated work into a shared computation across
non-adjacent transformations; \textbf{(3) downstream-aware pruning:}
remove work that is provably irrelevant to every final output;
\textbf{(4) pipeline-aware work placement:} move expensive work after
data-reducing steps as long as final outputs are not impacted;
\textbf{(5) rewrite-materialization co-optimization:} decide when
rewrite benefits are outweighed by materialization cost, and when
materializing a rewritten result makes the rewrite worthwhile; and
\textbf{(6) frequency-aware optimization:} align computation frequency
with input-change rates and output-consumption rates, avoiding
refreshes that do not impact consumers.

Structurally, \dagsmith first analyzes the compiled SQL and the
dependency edges to identify regions likely to contain these
opportunities. It then leverages LLMs to reason over the relevant
subgraph and propose concrete refactorings, but deliberately separates
proposal, criticism, SQL generation, and data-backed equivalence
checking, so that unsafe or counterproductive rewrites are rejected
before adoption. Adopted candidates are not evaluated in isolation:
\dagsmith retunes persistence choices for each candidate using a
learned cost model and an iterative local linearization procedure,
then solves a global selection problem to choose a non-conflicting
subset of rewrites. Finally, frequency information from source update
rates and output demand reweights cost estimates, detects
over-scheduled work, and generates splitting candidates when
slow-changing logic is embedded in fast-refreshing transformations.
The result is a rewritten, output-equivalent SQL pipeline that remains
executable by the same SQL engine and orchestration framework, but is
significantly more performant.

\ph{Contributions}
In summary, this chapter makes the following contributions:
\begin{enumerate}[leftmargin=*, topsep=2pt, itemsep=1pt]
\item We formulate dependency-aware source-to-source rewriting for SQL
pipeline DAGs. We present the first holistic approach to this problem
(to the best of our knowledge) by identifying the pipeline-level
signals exposed by explicit dependencies and the rewriting
opportunities they enable: dependency-edge simplification, non-local
semantic reuse, downstream-aware pruning, pipeline-aware work
placement, rewrite-materialization co-optimization, and
frequency-aware optimization.
\item We present the design of \dagsmith, which combines structural
graph analysis, LLM-guided refactoring with separated proposal,
criticism, SQL generation, and data-backed equivalence checking,
learned cost modeling, persistence tuning, frequency-aware scheduling,
and global candidate selection.
\item We conduct a comprehensive evaluation showing that, on the
open-source Tuva dbt project, \dagsmith reduces elapsed time by 42.6\%
and warehouse compute cost by 67.7\%; these reductions are
22.8\%/83.0\% larger than materialization-only tuning and
98.1\%/348.3\% larger than state-of-the-art single-query rewriting
techniques.
\end{enumerate}

The rest of the chapter proceeds as follows. \Cref{dag:sec:motivating}
defines the setting and illustrates the opportunity space,
\cref{dag:sec:approach} presents the design of \dagsmith, and
\cref{dag:sec:eval} presents our evaluation. Related work is discussed
in \cref{chap:related}, and \cref{dag:sec:conclusion} concludes the
chapter.
